% TAB_SUPERHEAVY_PREDICTIONS.tex
% Generated by: src/derivations/code/superheavy_predictions.py
% Date: 2026-02-01, Updated with p=6.0 (optimal prefactor)
% Model: V7.8 M2 (GN + d(n) coordination frustration)
% Audit: audit/radioactivity_v7_8_Salpha_deformation/07_FIT_RESULTS_V7_8.md
%
% MAJOR UPDATE: Prefactor changed from 6.1 to 6.0 based on sensitivity analysis.
% Result: Mean Δlog improved from 1.16 to 0.38 dex (3× better), 100% pass rate.

\begin{table}[htbp]
\centering
\caption{Superheavy $\alpha$-decay predictions using coordination-frustration-corrected
Geiger-Nuttall law [\textbf{I}]. Model extrapolated from $A \leq 257$ fit (V7.8 M2) to
$Z \geq 114$ regime. The frustration term $g \times d(n)$ with $g = -1.64 \pm 0.14$
dramatically improves predictions: baseline GN predicts $t_{1/2} \sim 10^6$~s (weeks)
for Og-294; frustration-corrected model predicts $\sim 0.9$~ms (experimental: 0.7~ms).
\textbf{Updated}: Prefactor optimized to $p = 6.0$ based on sensitivity analysis
(R$^2$ = 0.975, RMSE = 0.93).}
\label{tab:superheavy_predictions}
\small
\begin{tabular}{@{}llcccccccccc@{}}
\toprule
El. & $Z$ & $A$ & $n(A)$ & $d(n)$ & $Q_\alpha$ & $\log_{10} t$ & $\log_{10} t$ & $\pm\sigma$ & $t_\mathrm{pred}$ & $t_\mathrm{exp}$ & $\Delta$ \\
    &     &     &        &        & (MeV)      & (GN)          & (GN+$d$)      &             &                   &                  & log \\
\midrule
Fl  & 114 & 289 & 39.67 & 3.67 &  9.82 &  6.2 &  0.1   & 1.5 & 1.4 s   & 2.1 s      & \textbf{0.19} \\
Fl  & 114 & 290 & 39.71 & 3.71 &  9.19 &  8.1 &  2.0   & 1.6 & 1.6 min & $\sim$19 s & 0.70 \\
Mc  & 115 & 290 & 39.71 & 3.71 & 10.41 &  5.0 & $-1.1$ & 1.5 & 83 ms   & 0.65 s     & 0.89 \\
Lv  & 116 & 293 & 39.85 & 3.85 & 10.67 &  4.8 & $-1.5$ & 1.5 & 32 ms   & 60 ms      & \textbf{0.28} \\
Ts  & 117 & 294 & 39.90 & 3.90 & 10.81 &  5.0 & $-1.4$ & 1.5 & 35 ms   & 26 ms      & \textbf{0.13} \\
Og  & 118 & 294 & 39.90 & 3.90 & 11.65 &  3.4 & $-3.0$ & 1.5 & 0.9 ms  & 0.7 ms     & \textbf{0.12} \\
\midrule
\multicolumn{12}{c}{\textit{Predictions for unsynthesized elements}} \\
\midrule
119 & 119 & 298 & 40.08 & 4.08 & 10.5  &  6.8 &  0.1   & 1.7 & 1.1 s   & N/A     & --- \\
120 & 120 & 302 & 40.26 & 4.26 & 10.0  &  8.7 &  1.7   & 1.8 & 48 s    & N/A     & --- \\
120 & 120 & 304 & 40.34 & 4.34 &  9.5  & 10.2 &  3.1   & 2.0 & 20 min  & N/A     & --- \\
\bottomrule
\end{tabular}

\vspace{3mm}
\footnotesize
\textbf{Legend}:
$n(A) = 6.0 \times A^{1/3}$ [\textbf{Cal}] (optimized from 6.1);
$d(n) = \min_k |n(A) - 2^a 3^b|$ (nearest allowed coordination);
Forbidden zone: $n \in [37, 47]$ (all 11 integers);
$Q_\alpha$ from WS4+RBF/FRDM + exp. corrections.

\vspace{1mm}
\textbf{Pass criterion}: $|\Delta\log_{10}| < 1.5$ dex (factor $\sim$30).
Result: \textbf{6/6 pass (100\%)}.

\vspace{1mm}
\textbf{Mean improvement}: Baseline GN $\langle|\Delta|\rangle = 6.16$ dex $\to$
GN+$d(n)$ $\langle|\Delta|\rangle = 0.38$ dex (\textbf{16$\times$ better}).

\vspace{1mm}
\textbf{Data sources}: NUBASE2020, AME2020, Dubna SHE Factory (2024--2025),
Berkeley 88-Inch Cyclotron (2025), GSI/FAIR reports.

\vspace{1mm}
\textbf{Prefactor optimization}: Sensitivity analysis on 86 H0 nuclides (A $\leq$ 257)
shows $p = 6.0$ yields R$^2 = 0.975$, RMSE $= 0.93$ vs $p = 6.1$ with
R$^2 = 0.939$, RMSE $= 1.46$. See \texttt{prefactor\_sensitivity\_full.py}.
\end{table}
